\lecture{32}{Распределённая задача экономического диспетчирования}

\sourcevideo
    {https://www.youtube.com/playlist?list=PLOzRYVm0a65fhKDS8cYIC7YCuVGJ4FOJx}
    {Week 9: Lecture 32: Distributed Economic Dispatch Problem}

Экономическое диспетчирование состоит в распределении заданной
нагрузки между генераторами с минимальной суммарной стоимостью. Каждый
генератор знает собственную функцию затрат и обменивается с соседями
только локальными переменными. В этой лекции выводятся условия
оптимальности и формулируются требования к распределённому алгоритму;
конкретная динамика строится в следующей лекции.

\section{Постановка задачи}

Пусть в сети имеется \(N\) генераторов. Через \(P_i\) обозначим
мощность генератора \(i\), а через \(P_{\mathrm D}\) --- суммарную
нагрузку. В модели без потерь должен выполняться баланс
\[
    \sum_{i=1}^{N}P_i=P_{\mathrm D}.
\]
Генератору \(i\) соответствует функция стоимости \(C_i(P_i)\). В
лекции используется квадратичная модель
\[
    C_i(P_i)
    =
    \alpha_iP_i^2+\beta_iP_i+\gamma_i,
    \qquad
    \alpha_i>0,
\]
для которой
\[
    C_i'(P_i)=2\alpha_iP_i+\beta_i.
\]
Производная \(C_i'(P_i)\) называется предельной стоимостью генерации.
Постоянный коэффициент \(\gamma_i\) меняет значение стоимости, но не
оптимальное распределение мощности.

Полная централизованная задача имеет вид
\[
\begin{aligned}
    \min_{P_1,\ldots,P_N}\quad
    &\sum_{i=1}^{N}C_i(P_i),
    \\
    \text{при условиях}\quad
    &\sum_{i=1}^{N}P_i=P_{\mathrm D},
    \\
    &P_i^{\min}\le P_i\le P_i^{\max},
    \qquad i=1,\ldots,N.
\end{aligned}
\]
Необходимое условие допустимости равно
\[
    \sum_{i=1}^{N}P_i^{\min}
    \le
    P_{\mathrm D}
    \le
    \sum_{i=1}^{N}P_i^{\max}.
\]

Сначала ограничения мощности отбрасываются. Полученная задача
называется некапацитированной:
\[
\begin{aligned}
    \min_{P_1,\ldots,P_N}\quad
    &\sum_{i=1}^{N}C_i(P_i),
    \\
    \text{при условии}\quad
    &\sum_{i=1}^{N}P_i=P_{\mathrm D}.
\end{aligned}
\]
При строгой выпуклости функций \(C_i\) её оптимальный вектор
единственен.

В отличие от консенсусной оптимизации, здесь обычно
\[
    P_i^\star\ne P_j^\star.
\]
Согласовываются не мощности, а двойственная переменная, связанная с
общим балансом.

\section{Условия оптимальности}

Для ограничения
\[
    P_{\mathrm D}-\sum_{i=1}^{N}P_i=0
\]
введём множитель \(\lambda\in\mathbb R\). Функция Лагранжа равна
\[
    \mathcal L(P,\lambda)
    =
    \sum_{i=1}^{N}C_i(P_i)
    +
    \lambda
    \left(
        P_{\mathrm D}-\sum_{i=1}^{N}P_i
    \right).
\]
Для некапацитированной выпуклой задачи условия стационарности дают
\[
    C_i'(P_i^\star)=\lambda^\star,
    \qquad i=1,\ldots,N,
\]
а производная по \(\lambda\) восстанавливает баланс
\[
    \sum_{i=1}^{N}P_i^\star=P_{\mathrm D}.
\]
Следовательно,
\[
    C_1'(P_1^\star)
    =
    \cdots
    =
    C_N'(P_N^\star)
    =
    \lambda^\star.
\]

Если у двух свободных генераторов предельные стоимости различны, то
перенос малой мощности от более дорогого генератора к более дешёвому
сохраняет баланс и уменьшает общую стоимость. Поэтому равенство
предельных стоимостей является необходимым условием оптимальности.

\section{Квадратичная модель}

Для квадратичных затрат условие стационарности принимает вид
\[
    2\alpha_iP_i^\star+\beta_i=\lambda^\star,
\]
откуда
\[
    P_i^\star
    =
    \frac{\lambda^\star-\beta_i}{2\alpha_i}.
\]
Подстановка в условие баланса даёт
\[
    \sum_{i=1}^{N}
    \frac{\lambda^\star-\beta_i}{2\alpha_i}
    =
    P_{\mathrm D},
\]
и поэтому
\[
    \lambda^\star
    =
    \frac{
        P_{\mathrm D}
        +
        \displaystyle\sum_{i=1}^{N}
        \frac{\beta_i}{2\alpha_i}
    }{
        \displaystyle\sum_{i=1}^{N}
        \frac1{2\alpha_i}
    }.
\]
Зная \(\lambda^\star\), каждый генератор вычисляет собственную
оптимальную мощность только по своим коэффициентам. Однако прямое
вычисление общего множителя требует централизованного сбора
\(P_{\mathrm D}\), всех \(\alpha_i\) и всех \(\beta_i\). Именно от
этого сбора отказывается распределённая постановка.

\section{Распределённая формулировка}

Обмен сообщениями описывается связным графом
\[
    \mathcal G=(\mathcal V,\mathcal E),
    \qquad
    \mathcal V=\{1,\ldots,N\}.
\]
Ребро \(\{i,j\}\in\mathcal E\) означает наличие канала связи между
генераторами \(i\) и \(j\). Коммуникационный граф может совпадать с
физической сетью, но это не требуется.

Вместо одного глобального множителя генератор \(i\) хранит локальную
оценку \(\lambda_i(t)\). Целевое состояние удовлетворяет
\[
    \lambda_1(t)
    =
    \cdots
    =
    \lambda_N(t)
    \longrightarrow
    \lambda^\star,
\]
а текущая мощность связывается с ней формулой
\[
    P_i(t)
    =
    \frac{\lambda_i(t)-\beta_i}{2\alpha_i}.
\]
При \(\lambda_i(t)\to\lambda^\star\) получаем
\[
    P_i(t)\to P_i^\star.
\]

Обычного среднего консенсуса по \(\lambda_i\) недостаточно: его
предел определяется начальными значениями и не обязан совпадать с
\(\lambda^\star\). Распределённая динамика должна одновременно
согласовывать предельные стоимости и обеспечивать баланс мощности.

Обмен локальными двойственными оценками не требует явной передачи
коэффициентов \(\alpha_i,\beta_i,\gamma_i\), но сам по себе не даёт
формальной гарантии конфиденциальности.

\section{Сохранение баланса и ограничения мощности}

Баланс удобно обеспечить уже при инициализации. Пусть \(D_i\) ---
нагрузка, локально закреплённая за генератором \(i\), причём
\[
    \sum_{i=1}^{N}D_i=P_{\mathrm D}.
\]
Выбор
\[
    P_i(0)=D_i
\]
даёт допустимое начальное состояние. Если алгоритм удовлетворяет
\[
    \frac{\mathrm d}{\mathrm dt}
    \sum_{i=1}^{N}P_i(t)=0,
\]
то
\[
    \sum_{i=1}^{N}P_i(t)=P_{\mathrm D}
\]
для всех \(t\). Полная нагрузка при этом не обязана быть известна
каждому генератору отдельно.

В капацитированной задаче дополнительно требуется
\[
    P_i^{\min}\le P_i\le P_i^{\max}.
\]
Формула
\[
    P_i^\star
    =
    \frac{\lambda^\star-\beta_i}{2\alpha_i}
\]
применима только к генераторам с неактивными ограничениями. Если это
значение выходит за допустимый интервал, генератор работает на
соответствующей границе. Поэтому активные ограничения изменяют
условия оптимальности и требуют отдельного алгоритма.

\section{Требования к алгоритму и границы модели}

Для некапацитированной задачи целевое состояние определяется двумя
условиями:
\[
    C_i'(P_i^\star)=\lambda^\star,
    \qquad i=1,\ldots,N,
\]
и
\[
    \sum_{i=1}^{N}P_i^\star=P_{\mathrm D}.
\]
Следующая лекция строит распределённую динамику, которая сохраняет
баланс и за фиксированное время согласовывает предельные стоимости.
Генератор использует собственную функцию стоимости, текущую мощность,
локальную двойственную оценку и сообщения соседей; полные функции
стоимости остальных генераторов ему не нужны.

Рассматриваемая модель не учитывает потери в линиях, ограничения их
пропускной способности, динамику частоты и напряжения, стоимость
запуска и остановки, невыпуклые режимы, неопределённость
возобновляемой генерации и изменение нагрузки во времени. Поэтому она
служит базовой постановкой для введения распределённых методов, а не
полной моделью энергосистемы.
